\section{Head Analisys}
\label{app:head_analysis}

% -------------------------------------------------------------
\subsection{Spatial Aggregation and Dataset Representations}
\label{sec:aggregation}
% -------------------------------------------------------------

Each raw head output $\mathbf{H}_{\ell,b,h}$ is a three-dimensional tensor
indexed by window, token position, and feature dimension. To obtain a single
fixed-size vector per audio sample, we mean-pool over all windows and token
positions:
\begin{equation}
    \mathbf{r}_{\ell,b,h}
    = \frac{1}{N_w^\ell\, M}
      \sum_{i=1}^{N_w^\ell}\sum_{j=1}^{M}
      \mathbf{H}_{\ell,b,h}[i,j,:]
    \;\in\; \mathbb{R}^{d_h}.
    \label{eq:spatial_pooling}
\end{equation}
This yields one vector $\mathbf{r}_{\ell,b,h}$ per audio sample,
per head $h$, at every block $(\ell,b)$ of the network. Stacking these
vectors across the $n$ samples in the dataset gives the matrix
\begin{equation}
    \mathbf{R}_{\ell,b,h} \in \mathbb{R}^{n \times d_h},
\end{equation}
which is the primary object of our analysis. The total number of such
matrices across the model is $H_{\mathrm{tot}} = 184$ (one per head per
block; see Table~\ref{tab:architecture}).

Since $d_h = 24$ is constant across all stages, the ambient dimension of
$\mathbf{R}_{\ell,b,h}$ is uniform throughout the network, enabling direct
cross-stage comparisons of all dimensionality metrics without correction for
a varying ceiling.

\paragraph{Dataset sampling.}
We extract representations from four audio classification benchmarks:
\begin{itemize}
    \item \textbf{ESC-50}~\cite{piczak2015dataset}: 50 environmental sound classes, 2{,}000 clips
          (5\,s, 44.1\,kHz).
    \item \textbf{TinySOL}~\cite{romani2018tinysol}: 14 orchestral instrument classes with varied
          articulations, 2{,}071 monophonic samples (1--16\,s, 44.1\,kHz).
    \item \textbf{VocalSound}~\cite{gong2022vocalsound}: 6 non-speech vocal categories, stratified
          subset of 1{,}200 samples.
    \item \textbf{IRMAS}~\cite{bosch2012comparison}: 11 musical instrument classes, annotated
          audio excerpts with predominant instrument.
\end{itemize}
Audio preprocessing follows the CLAP standard pipeline: 64-band log-mel spectrogram ($f_{\min} =
50$\,Hz, $f_{\max} = 8000$\,Hz, FFT window 1024, hop 320), padded or truncated to 7 seconds at
44.1\,kHz.

% -------------------------------------------------------------
\subsection{Intrinsic Dimensionality Analysis}
\label{sec:dimensionality_method}
% -------------------------------------------------------------

To characterize the effective complexity of the projected head representations
$\{\mathbf{r}_{\ell,b,h}^{(i)}\}_{i=1}^{n} \subset \mathbb{R}^{1024}$, we employ a battery of
linear and nonlinear dimensionality estimators.

\subsubsection{Linear Estimators}

\paragraph{PCA-based dimensionality.}
For each head, let $\mathbf{R}_{\ell,b,h} \in \mathbb{R}^{n \times d_h}$ stack all aggregated
representations. We compute the covariance matrix $\mathbf{C}_{\ell,b,h} =
\tfrac{1}{n-1}\mathbf{R}^\top\mathbf{R}$ and obtain ordered eigenvalues $\lambda_1 \geq \lambda_2
\geq \cdots$. Linear intrinsic dimensionality is:
\begin{equation}
    d_{\mathrm{PCA}}(\alpha)
    = \arg\min_{k}\left\{
        \frac{\sum_{i=1}^{k}\lambda_i}{\sum_{i}\lambda_i} \geq \alpha
      \right\},
    \label{eq:pca_dim}
\end{equation}
evaluated at $\alpha \in \{0.90, 0.95, 0.99\}$. We also report the \emph{Explained Variance Ratio}
of the first principal component, $\mathrm{EVR}_1 = \lambda_1 / \sum_i \lambda_i$.

\paragraph{Participation Ratio.}
\begin{equation}
    \mathrm{PR}
    = \frac{\left(\sum_{i}\lambda_i\right)^2}{\sum_{i}\lambda_i^2}.
    \label{eq:pr}
\end{equation}
High PR signals uniform variance distribution; low PR signals dominance of few directions.

\paragraph{Effective Rank.}
\begin{equation}
    \mathrm{EffRank} = \exp\!\left(-\sum_{i} p_i \log p_i\right),
    \quad p_i = \frac{\lambda_i}{\sum_j \lambda_j},
    \label{eq:effrank}
\end{equation}
which can be interpreted as the exponential of the spectral entropy. Unlike PCA thresholds, it
provides a smooth, threshold-free notion of dimensionality, capturing how uniformly variance is
distributed across components.

\subsubsection{Nonlinear Estimators}

\paragraph{TwoNN.}
\begin{equation}
    d_{\mathrm{TwoNN}}
    = \left(\frac{1}{n}\sum_{i=1}^{n}\log\frac{r_2^{(i)}}{r_1^{(i)}}\right)^{-1},
    \label{eq:twonn}
\end{equation}
where $r_1^{(i)}, r_2^{(i)}$ are Euclidean distances to the first and second nearest neighbours of
$\mathbf{r}_{\ell,b,h}^{(i)}$~\cite{facco2017twonearest}.

\paragraph{MLE.}
\begin{equation}
    \hat{d}_{\mathrm{MLE}}(\mathbf{r})
    = \left(\frac{1}{k-1}\sum_{j=1}^{k-1}\log\frac{r_k(\mathbf{r})}{r_j(\mathbf{r})}\right)^{-1},
    \label{eq:mle}
\end{equation}
averaged over all samples~\cite{levina2005mle}, with $k=20$. This estimator is based on a
local maximum likelihood principle, assuming uniform sampling in small neighbourhoods, and is
sensitive to fine-scale manifold structure beyond global linear approximations.

\paragraph{Linear-Nonlinear (L/N) Ratio}
\begin{equation}
    \mathrm{Ratio}_B
    = \frac{\bar{d}_{\mathrm{PCA}_{99}}}{\bar{d}_{\mathrm{TwoNN}}}.
    \label{eq:ln_ratio}
\end{equation}
Values near 1 indicate near-linear manifolds; higher values signal nonlinear curvature beyond what
PCA captures, and serve as a diagnostic for selecting layer targets for spectral reweighting.